% !TeX root = ../../main.tex

Aufgrund der beschriebenen Grenzen von Retry-Mechanismen ist ein unkontrollierter Einsatz nicht ausreichend, um eine
zuverlässige Fehlerbehandlung zu gewährleisten.
Ein Retry-Mechanismus muss verschiedene Anforderungen erfüllen, um die Vorteile eines Retries nutzen zu können, ohne
dabei die Leistung und Stabilität der Datenbank negativ zu beeinflussen.

Eine wesentliche Anforderung ist die Zuverlässigkeit.
Der Retry-Mechanismus soll die Erfolgswahrscheinlichkeit fehlgeschlagener Transaktionen durch geeignete Retries
erhöhen, ohne dauerhaft fehlerhafte Transaktionen unnötig wiederholt auszuführen.
Dadurch kann unnötige Last auf die Datenbank vermieden und dauerhafte Fehler schneller erkannt werden.

Neben der Zuverlässigkeit spielt die Performance eine entscheidende Rolle.
Da jeder Wiederholungsversuch Ressourcen der Datenbank verbraucht, muss die Anzahl der Retries begrenzt und deren
Ausführung kontrolliert werden.
Die Belastung der Datenbank wird maßgeblich von Faktoren wie die maximale Anzahl an Wiederholungsversuchen oder
Wartezeiten zwischen dien einzelnen Versuchen beeinflusst.

Darüber hinaus sollte ein Retry-Mechanismus flexibel erweiterbar sein.
Unterschiedliche Fehlerarten benötigen unterschiedliche Behandlungsmethoden, da nicht jeder Fehler durch eine erneute
Ausführung behoben werden kann.
Eine optimierte Strategie muss daher an verschiedene Szenarien angepasst werden können.

Diese Anforderungen bilden die Grundlage für die Bewertung von Retry-Mechanismen in dieser Arbeit.
Dazu verdeutlichen sie, dass ein Retry-Mechanismus nicht ausschließlich auf eine erfolgreiche Wiederholung ausgelegt
sein darf.
Ebenso wichtig ist es, die Auswirkungen auf das Gesamtsystem zu berücksichtigen und einen Ausgleich zwischen
Fehlertoleranz und Systembelastung zu schaffen.
Nur durch eine geeignete Konfiguration kann ein Retry-Mechanismus langfristig zu einer stabilen udn performanten
Datenbankanwendung beitragen.
Da sich insbesondere die Anzahl der Wiederholungsversuche und die Wartezeit zwischen einzelnen Retries unmittelbar
auf die Belastung der Datenbank auswirken, werden diese Parameter im experimentellen Teil dieser Arbeit
systematisch untersucht.